\documentclass[a4paper,11pt]{article}
\usepackage[utf8]{inputenc}
\usepackage[T1]{fontenc}
\usepackage[french]{babel}
\usepackage[left=1cm,right=1cm,top=1cm,bottom=1cm]{geometry}
\usepackage{xcolor}
\usepackage{hyperref}
\usepackage{fontawesome5}
\usepackage{parskip}

% --- Couleurs GE HealthCare ---
\definecolor{primary}{RGB}{80, 35, 127}
\hypersetup{colorlinks=true, urlcolor=primary}

\begin{document}

% --- EXPÉDITEUR ---
\noindent
\begin{minipage}[t]{0.5\textwidth}
    \textbf{Loïc Aron MBASSI EWOLO}\\
    Élève-Ingénieur Bac+5 -- Double diplôme ENSEA\\[3pt]
    \faMapMarker\ Cergy, France\\
    \faPhone\ (+33) 7 59 52 33 98\\
    \faEnvelope\ \href{mailto:loicaron.mbassiewolo@ensea.fr}{loicaron.mbassiewolo@ensea.fr}
\end{minipage}
\hfill
% --- DATE ---
\begin{minipage}[t]{0.4\textwidth}
    \raggedleft
    Cergy, le 2 mars 2026
\end{minipage}

\vspace{1.2cm}

% --- DESTINATAIRE ---
\hfill
\begin{minipage}[t]{0.5\textwidth}
    \raggedleft
    \textbf{GE HealthCare France}\\
    Service Recrutement R\&D\\
    283 rue de la Minière, Buc (78)
\end{minipage}

\vspace{1cm}

\noindent\textbf{Objet :} Candidature en alternance R\&D -- IA pour la Santé / Maths Appliquées \& Traitement d'Images

\vspace{0.8cm}

Madame, Monsieur,

Classifier des anomalies cardiaques à partir d'images d'ECG avec TensorFlow, puis présenter les résultats à des mentors Google DeepMind : c'est le projet UROP que j'ai mené en 2024 et qui m'a orienté vers l'imagerie médicale. Ce week-end, le hackathon UNBOXED -- Agentic AI Medicine organisé à Centrale Lyon avec le soutien de GE HealthCare a confirmé cette orientation. L'échange direct avec vos mentors du site de Buc et avec des radiologues m'a donné une vision concrète des enjeux cliniques auxquels vos équipes R\&D répondent. Je souhaite prolonger cet engagement en rejoignant GE HealthCare en alternance, sur les sujets d'IA pour la santé ou de mathématiques appliquées et traitement d'images. Je tiens à votre disposition les travaux réalisés pendant le hackathon.

Mon profil technique est ancré dans l'imagerie médicale et le Deep Learning appliqué à la santé. Le projet UROP (classification d'anomalies ECG à partir d'images papier) m'a confronté au pipeline complet : étude bibliographique, extraction de features par traitement du signal (filtrage, analyse fréquentielle), conception du modèle de classification avec TensorFlow/Keras, et évaluation quantitative des performances. En parallèle, mes projets en Computer Vision (reconnaissance gestuelle avec PyTorch/OpenCV/Mediapipe, classification vidéo temps réel) m'ont formé aux problématiques de segmentation, détection et classification d'images 2D. Ma formation à l'ENSEA en traitement du signal et à Polytechnique Yaoundé en mathématiques appliquées (algèbre linéaire, optimisation, analyse numérique) couvre les fondamentaux théoriques nécessaires pour les thématiques de reconstruction 3D et de modélisation physique.

Le volet recherche de mon parcours correspond aux méthodes de travail d'un environnement R\&D : au MILA (Institut Québécois d'IA), j'ai mené une étude bibliographique SOTA, reproduit des expériences en PyTorch et rédigé des rapports de synthèse présentés à l'équipe. Le stage au Labo ICS de Florence sur l'analyse de signaux (Python, MATLAB) renforce cette expérience de recherche appliquée. Mon 1\textsuperscript{er} prix au hackathon IndabaX Africa 2025 (IA Santé : nettoyage de données médicales, déploiement d'API, dashboard) illustre ma capacité à livrer rapidement des prototypes fonctionnels dans le domaine médical.

Je suis disponible dès septembre 2026 pour une alternance sur votre site de Buc. La proximité avec Cergy facilite les déplacements quotidiens. L'opportunité de contribuer au développement d'algorithmes d'imagerie médicale au sein de GE HealthCare, en lien direct avec les experts cliniques, est la continuation logique de mon parcours.

Je reste à votre disposition pour un entretien afin de discuter de ma candidature.

Je vous prie d'agréer, Madame, Monsieur, l'expression de mes salutations distinguées.

\vspace{0.8cm}

\textbf{Loïc Aron MBASSI EWOLO}\\
\textit{Élève-Ingénieur ENSEA / Polytechnique Yaoundé}

\end{document}
